\section{问题四：限值由 10 降至 5 的电耗增幅}

\subsection{先做可行性判断}

在与问题二完全相同的历史电压和振打周期边界内重新求解 $5\mgNm$ 问题，R1、R7 和 R8 未通过独立 95\% 机会约束复核。它们分别对应高温异常、较高质量负荷和最高质量负荷。因此，若坚持历史边界，问题四不存在覆盖所有工况的统一增幅百分比；直接报告一个数会掩盖设备能力不足。

为评估进一步改造或提升整流变压器余量后的能耗代价，定义\textbf{条件情景}：四场历史最大电压上界统一放宽至 1.03 倍，振打周期上下界保持不变；其余工况、先验和随机检验完全相同。3\% 不是现场可直接执行的建议，而是局部能力扩展的仿真包络，实施前必须校核二次电流、火花率、绝缘和极板状态。

\subsection{各工况电耗增幅}

条件情景的结果见表\ref{tab:q4} 和图\ref{fig:q4}。数据字段 $P_{\mathrm{total}}$ 的单位为 kW，故表中直接比较的是总功率；在相同工况驻留时间下，电量 $E=P\Delta t$ 的增幅与功率增幅相同。8 类工况均通过 95\% 机会约束，等时电耗增幅为 5.22\%--6.11\%。

\begin{table}[H]
\centering
\caption{$10\rightarrow5\mgNm$ 的可行性、条件功率与等时电耗增幅}
\label{tab:q4}
\begin{tabular}{ccccccc}
\toprule
工况 & 占比/\% & $P_{10}$/kW & 历史边界下 $5$ 可行 & 条件 $P_5$/kW & 条件可行 & 增幅/\%\\
\midrule
R1 & 1.19 & 2049.8 & 否 & 2166.8 & 是 & 5.710\\
R2 & 16.90 & 1947.6 & 是 & 2066.5 & 是 & 6.105\\
R3 & 15.65 & 1989.7 & 是 & 2100.0 & 是 & 5.542\\
R4 & 16.02 & 2002.1 & 是 & 2117.2 & 是 & 5.751\\
R5 & 12.64 & 2073.3 & 是 & 2197.6 & 是 & 5.994\\
R6 & 17.02 & 1998.8 & 是 & 2103.2 & 是 & 5.220\\
R7 & 11.93 & 2110.9 & 否 & 2229.0 & 是 & 5.597\\
R8 & 8.63 & 2138.2 & 否 & 2251.2 & 是 & 5.285\\
\bottomrule
\end{tabular}
\end{table}

\begin{figure}[H]
\centering
\includegraphics[width=0.90\textwidth]{fig11_q4_energy_increase.png}
\caption{各工况 $10\rightarrow5\mgNm$ 的条件电耗增幅}
\label{fig:q4}
\end{figure}

以工况历史占比 $\pi_r$ 加权：
\begin{align}
\bar P_{10}&=\sum_{r=1}^{8}\pi_rP_{r,10}=2024.68\kW,\\
\bar P_{5}^{\mathrm{cond}}&=\sum_{r=1}^{8}\pi_rP_{r,5}^{\mathrm{cond}}=2139.21\kW,\\
\Delta P_{\%}&=100\%\times\frac{\bar P_{5}^{\mathrm{cond}}-\bar P_{10}}{\bar P_{10}}
=\boxed{5.66\%}.
\end{align}

最高负荷 R8 的条件增幅为 $5.29\%$。对该工况执行 12 组独立场景重采样，两级限值均全部可行，增幅的 95\% 区间为
\begin{equation}
\boxed{[5.05\%,\ 6.55\%]}.
\end{equation}
该区间反映情景抽样与优化不确定性，不包含设备老化、未观测二次电流或先验整体失配。

\subsection{增量功率归因}

将式\eqref{eq:power} 在 $10$ 与条件 $5\mgNm$ 策略间逐项相减，按工况频率加权，得到表\ref{tab:attrib}。$U_4$ 承担约一半的增量，说明限值收紧后末级精细捕集成为主要杠杆；周期项净增量接近零，是因为其主要用于重新分配尘负荷风险，而非直接承担平均电耗增长。

\begin{table}[H]
\centering
\caption{全周期增量功率的分项归因}
\label{tab:attrib}
\begin{tabular}{cccccccccc}
\toprule
分项 & $U_1$ & $U_2$ & $U_3$ & $U_4$ & $T_1$ & $T_2$ & $T_3$ & $T_4$ & 合计\\
\midrule
增量/kW & 24.91 & 12.23 & 19.28 & 58.14 & $-0.60$ & 0.23 & 0.26 & 0.07 & 114.53\\
占比/\% & 21.75 & 10.68 & 16.83 & 50.76 & $-0.52$ & 0.20 & 0.23 & 0.06 & 100.00\\
\bottomrule
\end{tabular}
\end{table}

出现小幅负归因不代表“发电”，而是两个最优策略间的场间重分配：某一场降低的功率被其他场更高效的捕集所替代。

\subsection{高浓度工况应对建议}

\begin{enumerate}
  \item \textbf{入口负荷前馈。} 用 $C_{\mathrm{in}}Q$ 预测高负荷到达时间，在烟气进入电场前预置 $U_1$--$U_3$ 和相应周期，减少控制滞后。
  \item \textbf{末级保留余量。} 常态不把 $U_4$ 长期推满，严格限值或短时峰值时由末级快速补偿；同步监测火花率和二次电流。
  \item \textbf{四场振打错峰。} 将触发相位错开，禁止多场尘饼同时脱落；高负荷时优先缩短尘负荷约束活跃电场的周期。
  \item \textbf{双阈值滞回。} 高负荷进入阈值高于退出阈值，最短保持时间不低于一个主要振打周期，避免频繁切换。
  \item \textbf{能力校核与应急模式。} 长期保证 $5\mgNm$ 前应完成整流变压器、电极、绝缘、火花率和灰斗输灰能力校核；若在线估计的 95\% 分位接近限值，则暂停非必要振打、提高末级电压并触发人工检查。
\end{enumerate}

\begin{warnbox}
\textbf{问题四最终口径：} 5.66\% 是“电压历史上界局部放宽 3\%，振打边界不变”的条件性全周期增幅。原始历史边界内有 3 类工况不可行，因此不存在无条件的全工况 $5\mgNm$ 最优增幅。
\end{warnbox}
